% Proposed B163, first saved33 source. No central edits.
\subsection*{The complete coset cohomology of the actual closed fold kernel}
\statement{B.163}{Proved}{The closed fold kernel has pointed-coset first cohomology, which injects into the actual deletion kernel and exactly measures its kernel on the split preimage}
\label{stmt:B163}
Use B.161--B.162 and take every closure inside the full completion of
the actual generic group. Put
\[
 \mathcal K_a=\overline{K_a},\qquad
 \mathcal P=\overline P=\mathcal K_a\rtimes
                   \widehat\lambda(\widehat J),\qquad
 \mathcal R=\overline R,\qquad \mathcal N=\overline N.
\]
For the continuous action of $\Delta=\langle\delta\rangle\cong C_2$,
there is a natural pointed homeomorphism
\begin{equation}\label{eq:B163-H1}
 \widehat J/\widehat S\ \xrightarrow{\ \sim\ }
 H^1(\Delta,\mathcal R),\qquad
 j\widehat S\longmapsto[X_{j\widehat S}].
\end{equation}
Here $H^1$ is the pointed space of cocycles $z\delta(z)=1$, modulo
$z\sim b^{-1}z\delta(b)$; it is not asserted to be a group.
The composite to $H^1(\Delta,\mathcal K_a)$ is injective, and its
image is exactly
\begin{equation}\label{eq:B163-image}
 \ker\bigl[H^1(\Delta,\mathcal K_a)
                 \longrightarrow H^1(\Delta,\mathcal P)\bigr].
\end{equation}
Moreover $\mathcal R\le\mathcal N\le\mathcal K_a$. Consequently
the map through $H^1(\Delta,\mathcal N)$ is injective as well, and
its image meets the kernel of
$H^1(\Delta,\mathcal N)\to H^1(\Delta,\mathcal K_a)$ only at the
basepoint.
\begin{proof}
\emph{Finite-factor torsion and compactness suffice.}
Use the explicit inverse system $U_V$ of B.162. At every finite coset
level, the semidirect product with the basis-inverting involution is
\[
 U_V\rtimes\Delta=\coprod_{x\in X_V}C_2(t_x),\qquad
 t_*=\tau,\quad t_x=X_x\tau.
\]
Both presentation maps are immediate from $(X_x\tau)^2=1$; this is
a finite-factor free profinite product. The prime-order torsion theorem
of \citet[Theorem A\(^{\prime}\)]{A_HR1985} says that every involution
is conjugate into a factor. The factor label is unique, since distinct
factor involutions have distinct coordinates in the finite elementary
abelian abelianization. For an orientation-odd involution a conjugator
can be chosen even: multiply an odd conjugator on the right by its
target factor involution. Thus every cocycle class at this level has
exactly one label in $X_V$.

For a cocycle $z\in U$, the labels of its finite-level images are
compatible, so determine $x\in X$. The closed subsets of $U$ whose
projections conjugate $X_{x_V}\tau$ to $z_V\tau$ are nonempty, because
$U\to U_V$ is onto. A common refinement handles any finite list of
these conditions. Compactness gives one even conjugator in $U$ for
all levels. Conversely different points of $X$ differ at a finite
level and cannot give cohomologous cocycles. This proves bijectivity
in \eqref{eq:B163-H1}. The map is continuous. The cocycle space is
compact and its compact-group orbit equivalence is closed, so $H^1$
is Hausdorff; the bijection is therefore a homeomorphism. Only the
finite-factor torsion theorem was used, not an infinite-index Kurosh
or abstract-kernel-completion identification.

\emph{Actual generic centralizers determine injectivity.}
In the full completion of the actual generic ordinary group, the fixed
subgroup is exactly $\widehat S$, by
B.123/B.159: in the orbifold free-product fiber the selected cone has
centralizer $\langle\tau\rangle$, also in full completion, and all
source operators fix it. In particular
\[
 \mathcal K_a^\delta=1,
\]
because the filled $q_a$ is injective on $\widehat S$.
The cocycle identity gives
\[
 X_{j\widehat S}\tau
   =\delta\widehat\lambda(j)\,\tau\,
                   \delta\widehat\lambda(j)^{-1}.
\]
If two such involutions are conjugate by $b\in\mathcal K_a$, then
\[
 \delta\widehat\lambda(j)^{-1}b\,
                   \delta\widehat\lambda(l)
\]
centralizes $\tau$, and hence belongs to $\widehat S$. Applying
$\widehat q_a$ gives $j^{-1}l\in\widehat S$. The labels coincide.
This proves injectivity into $H^1(\Delta,\mathcal K_a)$.

\emph{The whole kernel on the split preimage.}
Every $X_j$ is a coboundary in $\mathcal P$. Conversely suppose
$c\in\mathcal K_a$ is a cocycle and $c=p^{-1}\delta(p)$ for
$p\in\mathcal P$. Set $l=\widehat q_a(p)^{-1}$ and write
\[
 p=\delta\widehat\lambda(l)^{-1}k,\qquad k\in\mathcal K_a.
\]
This is possible since the filled target fixes $\widehat J$ pointwise.
Then
\[
 c=k^{-1}\delta\widehat\lambda(l)\widehat\lambda(l)^{-1}
                   \delta(k)=k^{-1}X_{l\widehat S}\delta(k).
\]
This proves \eqref{eq:B163-image}, retaining the order of every factor.

Finally the \emph{filled} factorization of B.139 gives $N\le K_a$.
B.161 gives $R\le N$, so the same inclusions hold for their actual
closures. Injectivity of the composite through $H^1(\Delta,\mathcal N)$
implies injectivity of its first map and the stated intersection with
the next kernel. It does not imply injectivity on all of
$H^1(\Delta,\mathcal N)$.
\end{proof}
The exact source distinction is bound in
\path{work_CE/pre_node_input_binding33.txt}.
The Herfort--Ribes input is the previously audited finite-factor
prime-order theorem; this proof adds only the displayed compact
finite-intersection argument.
